\section{Фактор-группы}


\subsection*{Классы смежности}

\begin{definition}
    \emph{Левое действие} подгруппы $H$ на $G$:
    \[
        H \to \mathrm{Bij}(G), \qquad h \mapsto \varphi_h,\quad \varphi_h(g) \coloneqq hg.
    \]
    Орбита элемента $g \in G$ при этом действии называется \emph{правым классом смежности}:
    \[
        \mathrm{Orb}\, g = \{hg \mid h \in H\} =: Hg.
    \]
    Множество всех правых классов смежности обозначается $H \backslash G$.
\end{definition}

\begin{remark}
    Действие здесь левое (умножаем на $h$ слева), но класс называется \emph{правым}, так как элемент
    $g$ стоит \emph{справа}.
\end{remark}

\begin{definition}
    \emph{Правое действие} подгруппы $H$ на $G$:
    \[
        H \to \mathrm{Bij}(G), \qquad h \mapsto \varphi_h,\quad \varphi_h(g) \coloneqq gh^{-1}.
    \]
    Орбита элемента $g$:
    \[
        \mathrm{Orb}\, g = \{gh^{-1} \mid h \in H\} = \{gh \mid h \in H\} =: gH
    \]
    называется \emph{левым классом смежности} элемента $g$.
    Множество всех левых классов смежности обозначается $G / H$.
\end{definition}

\begin{remark}
    Во втором случае необходимо писать именно $gh^{-1}$, а не $gh$. Если определить $\varphi_h(g) = gh$, то
    \[
        \varphi_{h_1 h_2}(g) = g h_1 h_2, \qquad (\varphi_{h_1} \circ \varphi_{h_2})(g) = \varphi_{h_1}
        (g h_2) = g h_2 h_1,
    \]
    и эти выражения не равны в общем случае. Обратные разворачивают порядок, и равенство восстанавливается:
    \[
        \varphi_{h_1 h_2}(g) = g(h_1 h_2)^{-1} = g h_2^{-1} h_1^{-1}
        = \varphi_{h_1}(g h_2^{-1}) = (\varphi_{h_1} \circ \varphi_{h_2})(g).
    \]
\end{remark}

\begin{remark}
    Чтобы не путать действие с умножением в группе, будем обозначать действие точкой: $h . g$
    (вместо формального $\varphi_h(g)$ или просто $hg$).
\end{remark}

% ----------------------------------------------------------------
\subsection*{Отношения эквивалентности}

\begin{stmt}
    Пусть $H < G$. Определим на $G$ два отношения:
    \[
        g_1 \sim_L g_2 \iff g_1,\, g_2 \in Hg \text{ для некоторого } g \in G,
    \]
    \[
        g_1 \sim_R g_2 \iff g_1,\, g_2 \in gH \text{ для некоторого } g \in G.
    \]
    Оба являются отношениями эквивалентности.

    \begin{proof}
        Докажем для $\sim_L$; для $\sim_R$ аналогично.

        \textbf{Рефлексивность.} $g \in Hg$ (при $h = e$), поэтому $g \sim_L g$.

        \textbf{Симметричность.} Если $g_1 \sim_L g_2$, то $g_1, g_2 \in Hg$, а значит и $g_2, g_1 \in Hg$,
        то есть $g_2 \sim_L g_1$.

        \textbf{Транзитивность.} Пусть $g_1 \sim_L g_2$ и $g_2 \sim_L g_3$, то есть
        \[
            g_1,\, g_2 \in Hg, \qquad g_2,\, g_3 \in H\widetilde{g}
        \]
        для некоторых $g, \widetilde{g} \in G$. Покажем, что $g_1, g_2 \in Hg_2$.

        Из $g_1 \in Hg$ следует $g_1 = h_1 g$, а из $g_2 \in Hg$ следует $g_2 = h_2 g$, откуда $g = h_2^{-1} g_2$.
        Подставляя:
        \[
            g_1 = h_1 g = h_1 h_2^{-1} g_2, \quad h_1 h_2^{-1} \in H,
        \]
        значит $g_1 \in Hg_2$. Очевидно, $g_2 \in Hg_2$.

        Аналогично из $g_2, g_3 \in H\widetilde{g}$ получаем $g_2, g_3 \in Hg_2$.

        Таким образом, $g_1, g_3 \in Hg_2$, откуда $g_1 \sim_L g_3$.
    \end{proof}
\end{stmt}

% ----------------------------------------------------------------
\subsection*{Свободность действий}

\begin{stmt}
    Пусть $H < G$. Оба действия $H$ на $G$ (левое и правое) свободны.

    \begin{proof}
        Рассмотрим левое действие. Если $h . g = hg = g$, то домножив справа на $g^{-1}$, получаем $h = e$.
        Для правого действия: если $g \cdot h^{-1} = g$, то $h^{-1} = e$, откуда $h = e$.
    \end{proof}
\end{stmt}

Итого, из структуры подгруппы $H < G$ автоматически получаем два свободных действия.

% ----------------------------------------------------------------
\subsection*{Теорема Лагранжа}

\begin{stmt}
    Все классы смежности имеют одинаковую мощность:
    \[
        |Hg| = |H| \quad \forall\, g \in G.
    \]
    \begin{proof}
        По формуле орбит--стабилизаторов: $|\mathrm{Orb}\, g| \cdot |\mathrm{Stab}\, g| = |H|$.
        Так как действие свободно, $|\mathrm{Stab}\, g| = 1$, откуда $|Hg| = |H|$.
    \end{proof}
\end{stmt}

\begin{theorem}[Лагранжа]
    Пусть $G$ --- конечная группа, $H < G$. Тогда $|H|$ делит $|G|$.

    \begin{proof}
        Группа $G$ разбивается в дизъюнктное объединение классов смежности:
        \[
            |G| = \sum_{\text{классы смежности}} |Hg| = |H| \cdot |G / H|.
        \]
        Значит, $|G| = |H| \cdot [G : H]$, и $|H| \mid |G|$.
    \end{proof}
\end{theorem}

\begin{definition}
    \emph{Индекс} подгруппы $H$ в группе $G$:
    \[
        [G : H] \coloneqq |G / H|
    \]
    --- количество левых (или правых) классов смежности.
\end{definition}

\begin{cor}
    Пусть $G$ --- конечная группа, $g \in G$. Тогда $\mathrm{ord}(g) = min \{n \in \mathbb{N} | g^n = e\}$ делит
    $|G|$, или считается равной бесконечности.

    \begin{proof}
        Рассмотрим циклическую подгруппу $\langle g \rangle = \{e, g, g^2, \ldots\} < G$.
        Её мощность равна $\mathrm{ord}(g)$, и по теореме Лагранжа $\mathrm{ord}(g) \mid |G|$.
    \end{proof}
\end{cor}


\subsection*{Нормальные подгруппы}

\begin{definition}
    Подгруппа $H < G$ называется \emph{нормальной}, если
    \[
        gH = Hg \quad \forall\, g \in G.
    \]
    Обозначение: $H \trianglelefteq G$.

    Эквивалентные условия (упражнение):
    \[
        H \trianglelefteq G \iff g^{-1}Hg = H \;\forall g \in G \iff g^{-1}Hg \subseteq H \;\forall g \in G.
    \]
    На практике удобно проверять именно последнее включение.
\end{definition}

\begin{stmt}
    Пусть $f\colon G \to T$ --- гомоморфизм. Тогда $\ker f \trianglelefteq G$.

    \begin{proof}
        Возьмём произвольные $g \in G$, $h \in \ker f$. Проверим, что $g^{-1}hg \in \ker f$:
        \[
            f(g^{-1}hg) = f(g^{-1}) \cdot f(h) \cdot f(g) = f(g)^{-1} \cdot e \cdot f(g) = e. \qedhere
        \]
    \end{proof}
\end{stmt}

\begin{remark}
    Верно и обратное: любая нормальная подгруппа является ядром некоторого гомоморфизма. Мы увидим это из следующей теоремы.
\end{remark}



\subsection*{Фактор-группа}

Хотим определить операцию на $G/H$. Естественный кандидат:
\[
    (g_1 H) \cdot (g_2 H) \coloneqq (g_1 g_2) H. \tag{$*$}
\]
Проблема: представитель не единственен, нужна корректность.

\begin{stmt}
    Формула $(*)$ определяет структуру группы на $G/H$ тогда и только тогда, когда $H \trianglelefteq G$.

    \begin{proof}
        \textbf{($\Leftarrow$) Корректность.} Пусть $g_1, t_1 \in g_1 H = t_1 H$ и $g_2, t_2 \in g_2 H = t_2 H$.
        Покажем, что $t_1 t_2 H = g_1 g_2 H$.

        \[
        (g_1 H)\cdot(g_2 H) = (g_1 g_2)H
        \]

        \[
        (t_1 t_2)H = t_1(t_2 H) = t_1(g_2 H) =
        \]

        \[
        = t_1(H g_2) = (t_1 H) g_2 = (g_1 H) g_2 =
        \]

        \[
        = (g_1 g_2)H
        \]
        Значит, $t_1 t_2 H = g_1 g_2 H$. Операция корректно определена.

        \textbf{Аксиомы группы.}
        \begin{itemize}
            \item \emph{Ассоциативность} --- бесплатно из ассоциативности в $G$.
            \item \emph{Нейтральный элемент} --- $eH$.
            \item \emph{Обратный} --- $(gH)^{-1} = g^{-1}H$.
        \end{itemize}

        \textbf{($\Rightarrow$)} Рассмотрим отображение $\pi\colon G \to G/H$, $g \mapsto gH$. Если $(*)$ задаёт операцию группы,
        то $\pi$ --- гомоморфизм:
        \[
            \pi(g_1 g_2) = g_1 g_2 H = (g_1 H)(g_2 H) = \pi(g_1)\pi(g_2).
        \]
        Его ядро:
        \[
            \ker \pi = \{g \in G \mid gH = eH\} = \{g \in G \mid g \in H\} = H.
        \]
        Следовательно, $H = \ker \pi \trianglelefteq G$.
    \end{proof}
\end{stmt}

\begin{definition}
    Группа $G/H$ называется \emph{фактор-группой} $G$ по нормальной подгруппе $H \trianglelefteq G$.
\end{definition}

% ----------------------------------------------------------------
\subsection*{Универсальное свойство фактор-группы}

\begin{stmt}[Универсальное свойство фактор-группы]
    Пусть $H \trianglelefteq G$, $T$ --- группа, $f\colon G \to T$ --- гомоморфизм такой, что $f(H) = \{e\}$.
    Тогда существует единственный гомоморфизм $\bar{f}\colon G/H \to T$ такой, что $f = \bar{f} \circ \pi$:
    \[
        \begin{array}{ccc}
            G & \xrightarrow{\ f\ } & T \\
            \downarrow\, \pi & \nearrow\, \bar{f} & \\
            G/H & &
        \end{array}
    \]

    \begin{proof}
        \textbf{Определение.} Положим $\bar{f}(gH) \coloneqq f(g)$.

        \textbf{Корректность.} Если $g' \in gH$, то $g' = gh$ для некоторого $h \in H$, откуда
        \[
            f(g') = f(gh) = f(g) f(h) = f(g) \cdot e = f(g).
        \]

        \textbf{Гомоморфизм.} $\bar{f}((g_1 H)(g_2 H)) = \bar{f}(g_1 g_2 H) = f(g_1 g_2) = f(g_1) f(g_2) = \bar{f}(g_1 H)\bar{f}(g_2 H)$.

        \textbf{Коммутативность.} $(\bar{f} \circ \pi)(g) = \bar{f}(gH) = f(g)$.

        \textbf{Единственность.} Любой $\psi\colon G/H \to T$, решающий задачу, обязан удовлетворять
        $\psi(gH) = \psi(\pi(g)) = f(g)$. Значит, $\psi = \bar{f}$.
    \end{proof}
\end{stmt}

% ----------------------------------------------------------------
\subsection*{Коммутаторы и коммутант}

\begin{definition}
    \emph{Коммутатором} элементов $g, t \in G$ называется
    \[
        [g, t] \coloneqq g t g^{-1} t^{-1}.
    \]
\end{definition}

\begin{stmt}
    Группа $G$ абелева тогда и только тогда, когда $[g, t] = e$ для всех $g, t \in G$.

    \begin{proof}
        \textbf{($\Rightarrow$)} Если $G$ абелева, то $gtg^{-1}t^{-1} = gg^{-1}tt^{-1} = e$.

        \textbf{($\Leftarrow$)} Если $[g, t] = gtg^{-1}t^{-1} = e$, то $gt = tg$ для любых $g, t \in G$.
    \end{proof}
\end{stmt}

\begin{remark}
    Коммутаторы не образуют подгруппу: произведение двух коммутаторов не обязано само быть коммутатором.
\end{remark}

\begin{definition}
    \emph{Коммутант} группы $G$:
    \[
        [G, G] \coloneqq \langle [g, t] \mid g, t \in G \rangle
    \]
    --- наименьшая подгруппа, содержащая все коммутаторы. Элементы $[G, G]$ --- это конечные произведения коммутаторов.
\end{definition}

\begin{stmt}
    $[G, G] \trianglelefteq G$.

    \begin{proof}
        Достаточно показать, что $g [t_1, t_2] g^{-1} \in [G, G]$ для любых $g, t_1, t_2 \in G$ (тогда сопряжение произвольного произведения
        коммутаторов разбивается на сопряжения каждого множителя).

        Заметим, что:
        \[
            g [t_1, t_2] g^{-1} = [g, [t_1, t_2]] \cdot [t_1, t_2].
        \]
        Правая часть --- произведение двух элементов коммутанта, значит $g[t_1, t_2]g^{-1} \in [G,G]$.

        Для произведения $[t_1, t_2][t_3, t_4] \cdots$:
        \[
            g \bigl([t_1,t_2][t_3,t_4]\cdots\bigr) g^{-1}
            = (g[t_1,t_2]g^{-1})(g[t_3,t_4]g^{-1})\cdots \in [G, G]. \qedhere
        \]
    \end{proof}
\end{stmt}

\begin{definition}
    \emph{Абелианизация} группы $G$ --- это фактор-группа
    \[
        G^{\mathrm{ab}} \coloneqq G / [G, G].
    \]
    Это наибольшая абелева группа, которую можно получить из $G$ факторизацией.
\end{definition}

% ----------------------------------------------------------------
\subsection*{Универсальное свойство абелианизации}

\begin{stmt}[Универсальное свойство абелианизации]
    Пусть $G$ --- группа, $A$ --- абелева группа, $f\colon G \to A$ --- гомоморфизм.
    Тогда существует единственный гомоморфизм $\bar{f}\colon G/[G,G] \to A$ такой, что $f = \bar{f} \circ \pi$:
    \[
        \begin{array}{ccc}
            G & \xrightarrow{\ f\ } & A \\
            \downarrow\, \pi & \nearrow\, \bar{f} & \\
            G/[G,G] & &
        \end{array}
    \]

    \begin{proof}
        Это следствие универсального свойства фактор-группы. Достаточно проверить, что $f([G,G]) = \{e\}$.

        Так как $A$ абелева, для любого коммутатора $[g_1, g_2] = g_1 g_2 g_1^{-1} g_2^{-1}$:
        \[
            f([g_1, g_2]) = f(g_1 g_2 g_1^{-1} g_2^{-1}) = f(g_1) f(g_2) f(g_1^{-1}) f(g_2^{-1}) =
        \]
        \[
            = f(g_1) f(g_1^{-1}) f(g_2) f(g_2^{-1}) = f(g_1 g_1^{-1} g_2^{-1}) = f(e) = e.
        \]

        Значит, все коммутаторы отображаются в $e$, а значит и весь $[G,G]$.
        По универсальному свойству фактор-группы такой $\bar{f}$ существует и единственен.
    \end{proof}
\end{stmt}


\section{Симметрические группы}

\subsection*{Знак перестановки}

\begin{definition}
    $S_n = \mathrm{Bij}(\{1, \ldots, n\}).$
    Пусть $\sigma \in S_n$. \emph{Число инверсий} перестановки $\sigma$:
    \[
        \mathrm{inv}(\sigma) \coloneqq \bigl|\{(i,j) \mid i < j,\ \sigma(i) > \sigma(j)\}\bigr|.
    \]
    \emph{Знак} перестановки:
    \[
        \mathrm{sgn}(\sigma) \coloneqq (-1)^{\mathrm{inv}{\sigma}} \in \{1, -1\}.
    \]
    Перестановка называется \emph{чётной}, если $\mathrm{sgn}(\sigma) = 1$, и \emph{нечётной} --- иначе.
\end{definition}

\begin{definition}
    Транспозиция --- перестановка, которая меняет местами два элемента.
\end{definition}

\begin{example}
    Рассмотрим транспозицию $\sigma = (24) \in S_5$ (меняет $2$ и $4$ местами, остальные на месте):
    \[
        \sigma\colon 1 \mapsto 1,\quad 2 \mapsto 4,\quad 3 \mapsto 3,\quad 4 \mapsto 2,\quad 5 \mapsto 5.
    \]
    Инверсии возникают в парах $(i,j)$, $i < j$, где порядок нарушается. Перебирая:
    \begin{center}
    \begin{tabular}{c|c|c}
        $(i,j)$ & было & стало \\
        \hline
        $(2,3)$ & $2 < 3$ & $4 > 3$ --- инверсия \\
        $(2,4)$ & $2 < 4$ & $4 > 2$ --- инверсия \\
        $(3,4)$ & $3 < 4$ & $3 > 2$ --- инверсия \\
    \end{tabular}
    \end{center}
    Итого $\mathrm{inv}(\sigma) = 3$, поэтому $\mathrm{sgn}(\sigma) = (-1)^3 = -1$.
\end{example}

\begin{stmt}
    У любой транспозиции $\tau \in S_n$ знак равен $-1$.
\end{stmt}

\begin{stmt}
    Знак перестановки мультипликативен:
    \[
        \mathrm{sgn}(\sigma_1 \sigma_2) = \mathrm{sgn}(\sigma_1)\cdot\mathrm{sgn}(\sigma_2).
    \]
    Иными словами, $\mathrm{sgn}\colon S_n \to \{1,-1\}$ --- гомоморфизм групп.

    Отсюда немедленно следует таблица:
    \[
        \text{чётн} \cdot \text{чётн} = \text{чётн}, \qquad
        \text{нечётн} \cdot \text{нечётн} = \text{чётн}, \qquad
        \text{чётн} \cdot \text{нечётн} = \text{нечётн}.
    \]
\end{stmt}

\begin{stmt}
    У цикла длины $m$ знак равен $(-1)^{m-1}$.
\end{stmt}

\begin{example}
    Рассмотрим $\sigma = (234) \in S_5$ (цикл $2\to3\to4\to2$, элементы $1,5$ неподвижны):
    \[
        \sigma\colon 1\mapsto 1,\quad 2\mapsto 3,\quad 3\mapsto 4,\quad 4\mapsto 2,\quad 5\mapsto 5.
    \]

    \textbf{Способ 1: подсчёт инверсий.}
    Инверсии среди пар $(i,j)$, $i < j$: $(2,4)$ и $(3,4)$. Итого $\mathrm{inv}(\sigma) = 2$, $\mathrm{sgn}(\sigma) = 1$.

    \textbf{Способ 2: разложение в транспозиции.}
    $\sigma = (34)(24)$, два множителя:
    \[
        \mathrm{sgn}(\sigma) = \mathrm{sgn}((34) \cdot (24)) = \mathrm{sgn}(34)\cdot\mathrm{sgn}(24) = (-1)^2 = 1.
    \]

    \textbf{Способ 3: формула для цикла.}
    Цикл длины $3$: $\mathrm{sgn} = (-1)^{3-1} = 1$.
\end{example}



\begin{definition}
    $A_n \coloneqq \ker(\mathrm{sgn}) \subset S_n$ --- подгруппа группы $S_n$, состоящая из чётных перестановок.
\end{definition}

\begin{stmt}
    $A_n \trianglelefteq S_n$, $[S_n : A_n] = 2$, $|A_n| = n!/2$ (при $n \geq 2$).
    \begin{proof}
        Нормальность: $A_n = \ker(\mathrm{sgn})$ --- ядро гомоморфизма.

        Индекс: по следствию из теоремы о гомоморфизмах $|\mathrm{Im}(\mathrm{sgn})| = |S_n|/|A_n|$.
        При $n \geq 2$ образ равен $\{1,-1\}$ (транспозиция существует), значит $|S_n|/|A_n| = 2$.
    \end{proof}
\end{stmt}


\subsection*{Коммутант симметрической группы}

\begin{lemma}
    Знак любого коммутатора равен $1$.

    \begin{proof}
        Знак --- гомоморфизм в абелеву группу $\{1,-1\}$, поэтому:
        \[
            \mathrm{sgn}([\sigma,\tau])
            = \mathrm{sgn}(\sigma\tau\sigma^{-1}\tau^{-1})
            = \mathrm{sgn}(\sigma)\,\mathrm{sgn}(\tau)\,\mathrm{sgn}(\sigma)^{-1}\,\mathrm{sgn}(\tau)^{-1}
            = 1,
        \]
        так как в абелевой группе $\{1,-1\}$ можно переставлять множители.
    \end{proof}
\end{lemma}

\begin{cor}
    $[S_n, S_n] \subseteq A_n$.
\end{cor}

\begin{stmt}
    $[S_3, S_3] = A_3$.
    \begin{proof}
        Сначала установим, что $[S_3, S_3] \subseteq A_3$ (любой коммутатор чётен, по лемме выше).

        Далее, $|S_3| = 6$, $|A_3| = 3$, так что $A_3 = \{e, (123), (132)\}$ --- все 3-циклы.

        Так как $[S_3, S_3] \leq A_3$ и $|A_3| = 3$, по теореме Лагранжа $|[S_3, S_3]| \in \{1, 3\}$.

        Если $|[S_3, S_3]| = 1$, то $[S_3, S_3] = \{e\}$, то есть $S_3$ абелева --- противоречие (например,
        $(12)(23) \neq (23)(12)$).

        Значит $|[S_3, S_3]| = 3$, откуда $[S_3, S_3] = A_3$.
    \end{proof}
\end{stmt}

\begin{lemma}
    $A_n$ порождается циклами длины $3$.
    \begin{proof}
        Пусть $T \leq S_n$ --- подгруппа, порождённая всеми циклами длины $3$.

        \textbf{$T \subseteq A_n$.} Цикл длины $3$ чётен ($\mathrm{sgn} = (-1)^2 = 1$), значит все его произведения тоже чётны.

        \textbf{$A_n \subseteq T$.} Возьмём произвольную $\sigma \in A_n$. Запишем её как произведение транспозиций:
        \[
            \sigma = \tau_1 \tau_2 \cdots \tau_N.
        \]
        Так как $\mathrm{sgn}(\sigma) = 1 = (-1)^N$, число $N$ чётно. Разобьём транспозиции попарно и покажем,
        что произведение любых двух транспозиций лежит в $T$.

        Пусть $\tau_1 = (ab)$, $\tau_2 = (cd)$. Возможны три случая:

        \textbf{Случай 1: $(ab) = (cd)$.} Тогда $\tau_1\tau_2 = e \in T$.

        \textbf{Случай 2: $|\{a,b\}\cap\{c,d\}|=1$.} НУО $b = c, a \ne d$. Тогда:
        \[
            (ab)(bd) = (abd).
        \]
        Проверка:
        \[
            a \xrightarrow{(bd)} a \xrightarrow{(ab)} b, \quad
            b \xrightarrow{(bd)} d \xrightarrow{(ab)} d, \quad
            d \xrightarrow{(bd)} b \xrightarrow{(ab)} a.
        \]
        Таким образом $(ab)(bd) = (a\ b\ d)$ --- цикл длины $3$, лежит в $T$.

        \textbf{Случай 3: $\{a,b\}\cap\{c,d\}=\varnothing$.} Тогда:
        \[
            (ab)(cd) = (acb)(acd).
        \]
        Проверка ($a,b,c,d$ --- различные элементы, сначала $(cd)$, потом $(ab)$):
        \[
            a \xrightarrow{(cd)} a \xrightarrow{(ab)} b, \quad
            b \xrightarrow{(cd)} b \xrightarrow{(ab)} a, \quad
            c \xrightarrow{(cd)} d \xrightarrow{(ab)} d, \quad
            d \xrightarrow{(cd)} c \xrightarrow{(ab)} c.
        \]
        С другой стороны, $(acb)(acd)$:
        \[
            a \xrightarrow{(acd)} c \xrightarrow{(acb)} b, \quad
            b \xrightarrow{(acd)} b \xrightarrow{(acb)} a, \quad
            c \xrightarrow{(acd)} d \xrightarrow{(acb)} d, \quad
            d \xrightarrow{(acd)} a \xrightarrow{(acb)} c.
        \]
        Совпадает.  Произведение двух 3-циклов лежит в $T$.

        Итого, каждая пара транспозиций даёт элемент $T$, значит $\sigma \in T \Rightarrow A_n = T$.
    \end{proof}
\end{lemma}

\begin{theorem}
    $[S_n, S_n] = A_n$.
    \begin{proof}
        Уже знаем $[S_n, S_n] \subseteq A_n$.

        Для обратного включения: по утверждению выше для любых трёх элементов $\{i,j,k\} \subset \{1,\ldots,n\}$ группа
        $S_{\{i,j,k\}} \cong S_3$ вкладывается в $S_n$ как подгруппа, переставляющая только $i,j,k$. По утверждению ниже коммутант
        $S_3$ --- это $A_3$, состоящий ровно из двух 3-циклов на $\{i,j,k\}$. Значит, все 3-циклы лежат в
        $[S_n, S_n]$.

        По лемме $A_n$ порождается 3-циклами, откуда $A_n \subseteq [S_n, S_n]$.
    \end{proof}
\end{theorem}

